\chapter{The resonance problem}
\label{sec:introduction}
% BEGIN CURATED SUBJECT INDEX
\index{operator!adjoint and self-adjoint}
\index{Hilbert space}
\index{topology}
\index{measure theory}
\index{spectral theorem}
\index{spectrum!continuous}
\index{spectrum!point}
\index{resolvent}
\index{Riesz projection}
\index{direct integral}
% END CURATED SUBJECT INDEX

An unstable quantum state does not fit comfortably into the spectral theorem
for a self-adjoint Hamiltonian.  A normalized eigenvector of a self-adjoint
operator has a real eigenvalue and evolves only by a phase.  A state that loses
probability from a localized region instead has, over an intermediate time
window, an amplitude of the form
\begin{equation}
  A(t)\simeq Z\rme^{-\frac{\rmi E_* t}{\hbar}
  -\frac{\Gamma t}{2\hbar}},
  \qquad t>0 .
  \label{eq:intro-decay}
\end{equation}
Here $E_*$ is the resonance position, $\Gamma>0$ is the decay width, and $Z$
is the pole residue projected onto the prepared state.  It is tempting to call
$E_*-\rmi\Gamma/2$ an eigenvalue.  Yet the corresponding outgoing wave grows
exponentially in space and is not an element of the original Hilbert space.
The apparent contradiction is resolved only after one specifies which object
has been analytically continued, which function space is being used, and which
boundary value of the continuum resolvent is meant.

This issue is older than the modern language of non-Hermitian physics.
Siegert's outgoing states, Zel'dovich's Gaussian regularization, Berggren's
completeness relation, Feshbach's energy-dependent effective Hamiltonian, and
the Aguilar--Balslev--Combes construction all address different parts of the
same problem
\cite{Siegert:1939,Zel'dovich:1961,Berggren:1968zz,
Feshbach1958UnifiedI,Feshbach1962UnifiedII,Aguilar:1971ve,Balslev:1971vb}.
The language has broadened because non-self-adjoint generators now arise in
nuclear reactions, atomic and molecular ionization, optical resonators,
open-system dynamics, topological phases, and black-hole perturbation theory
\cite{Myo:2014ypa,Rotter:2007zng,Ashida2020NonHermitian,
Berti:2009kk}.  In all of these settings, complex eigenvalues are useful, but
they are not by themselves a definition of a resonance.

The resurgence program changes the order in which this problem can be
understood.  Exact WKB analysis starts from a formal semiclassical solution,
Borel resums it in sectors of the complex coordinate plane, and imposes global
consistency through Stokes automorphisms and monodromy.  Stable spectra,
tunneling splittings, and nonperturbative ambiguities are naturally expressed
by Voros symbols and exact quantization conditions
\cite{Voros:1983xx,Delabaere:1999xx,Iwaki:2014vad}.
The extension to resonances shows that an outgoing boundary condition can be
read as the vanishing of the prohibited incoming coefficient after analytic
continuation.  The complex energy then follows from the same connection
algebra that quantizes a bound state
\cite{Morikawa:2025grx,Morikawa:2025xjq,Morikawa:2025vvs}.
This observation is the organizing principle of the book.  It lets
us construct the analytic resonance condition first and introduce a
non-self-adjoint eigenvalue problem, a regularized norm, or a generalized
state only when one of those representations is needed for a subsequent
calculation.

Our purpose is therefore not to place several resonance formalisms side by
side.  We derive them in a dependency order.  The
primary calculational object is a global connection coefficient, or
equivalently the analytic determinant whose zero removes the incoming wave.
After continuation this zero is a pole of the resolvent and scattering
matrix, accompanied by a residue and a nearby continuum.  The resulting
dictionary is
\begin{align}
  &\text{zero of a Jost function}
  \quad\Longleftrightarrow\quad
  \text{pole of the scattering matrix}
  \notag\\
  &\quad\Longleftrightarrow\quad
  \text{outgoing Gamow--Siegert vector}
  \notag\\
  &\quad\Longleftrightarrow\quad
  \text{discrete eigenvalue of a complex-scaled operator}
  \notag\\
  &\quad\Longleftrightarrow\quad
  \text{zero of an exact-WKB connection coefficient} .
  \label{eq:intro-dictionary}
\end{align}
The last entry is placed last in the displayed dictionary only for typographic
convenience; in the derivation it comes first.  Each arrow has hypotheses.
Thresholds introduce branch points; long-range
potentials modify asymptotic states; complex scaling requires an analytic
deformation; a rigged Hilbert space depends on the test-function topology; and
an exact-WKB relation requires control of Borel summability and Stokes data.
The textbook aim is to make the reader verify these hypotheses during the
calculation, rather than encounter them later as qualifications to a formal
answer.

There is also a useful mathematical lesson behind this qualification.  A
Hilbert space alone does not specify which limits, generalized vectors, or
observables are admissible.  Norm, weak, strong-operator, and weak-operator
topologies answer different physical questions.  The spectral theorem is most
transparent as a direct integral over energy, with a fiber that records open
channels and spectral multiplicity.  The \vN\ algebra generated by the
spectral projections is the algebra of essentially bounded functions on this
measure space, while operators commuting with the Hamiltonian act fiber by
fiber.  In that representation the formal statement ``$S$ is diagonal in
energy'' becomes the precise decomposition
$\Scattering=\int^\oplus \Scattering(E)\dd\mu(E)$.
Resonances then appear exactly at the boundary of the self-adjoint theorem:
they are not additional spectral projections of $H$, but residues obtained
after analytic continuation or Riesz projections of a deformed operator.

Long-range forces expose a second boundary.  For $V(r)\sim 1/r$, the phase
accumulated along a free trajectory diverges logarithmically, so the ordinary
M{\o}ller comparison with a free plane wave fails.  Two-body Coulomb scattering
is repaired by a known logarithmic phase and Coulomb wave.  For three charged
particles there are several inequivalent asymptotic regions, and no single
elementary Coulomb plane wave is uniformly valid across all of them.  In QED,
the analogous logarithm is carried by infinitely many soft photons: a charged
asymptotic state generally belongs to a coherent, non-Fock representation of
the canonical commutation relations.  These facts motivate treating topology,
measure, representation, and asymptotic dynamics as physical data rather than
background formalism.

The same dictionary clarifies why black-hole quasinormal modes are quantum
resonances rather than merely damped normal modes.  With time dependence
$\rme^{-\rmi\omega t}$, a black-hole perturbation is ingoing at the event
horizon and outgoing at infinity or at a cosmological horizon.  The resulting
frequency $\omega$ lies in the lower half-plane.  The radial wave function is
not square integrable on the original tortoise-coordinate line, but becomes
normalizable after an admissible complex deformation.  Complex scaling then
exposes quasinormal poles and rotates the continuum in the variable
$E=\omega^2$
\cite{Ogawa:2026veu,Ogawa:2026dSCSM}.  This makes it possible to discuss pole
frequencies and continuum response in one spectral calculation.

Exceptional points provide a second extension.  When parameters are varied,
a resonance pole may meet another discrete pole or a state on the rotated
continuum.  At the coalescence, eigenvalues and eigenvectors merge and the
operator develops a Jordan chain.  Encircling the branch point exchanges
spectral sheets and produces a geometric holonomy.  In an unbounded scattering
problem, however, continuum normalization adds a regulator-dependent factor
that is absent from a finite matrix.  A momentum-bin construction and a
monodromy-renormalization condition isolate the physical holonomy~\cite{Morikawa:2025inx,Garmon2021AnomalousEP}.  This is a useful example of how
non-Hermitian few-level intuition must be amended, rather than discarded, in
an infinite-dimensional problem.

Figure~\ref{fig:unified-plane} summarizes the spectral geometry used
throughout the book.  The continuous spectrum on the positive real axis is
the boundary of different analytic sheets.  A complex deformation rotates
that spectrum, revealing poles that were hidden behind the cut.  Exact WKB
analysis reaches the same poles by continuing solutions on the complex
coordinate plane.  A rigged Hilbert space records the corresponding pole
functionals without pretending that they are ordinary $L^2$ vectors.

\begin{figure}[t]
  \centering
  \begin{tikzpicture}[x=1.05cm,y=1.0cm,>=Latex]
    \draw[->] (-3.2,0) -- (5.2,0) node[right] {$\re E$};
    \draw[->] (0,-3.0) -- (0,2.4) node[above] {$\im E$};
    \draw[very thick,teal!70!black] (0,0) -- (4.7,0)
      node[pos=0.77,above] {continuum cut};
    \draw[very thick,blue!65!black,->] (0,0) -- (3.8,-2.2)
      node[below right] {$\rme^{-2\rmi\theta}[0,\infty)$};
    \fill[black] (-1.6,0) circle (2.2pt)
      node[above] {bound state};
    \fill[red!75!black] (2.0,-0.9) circle (2.4pt)
      node[below left] {$E_*-\rmi\Gamma/2$};
    \draw[red!75!black,dashed] (2.0,-0.9) -- (2.0,0);
    \draw[decorate,decoration={snake,amplitude=0.7mm,segment length=3mm},
      orange!80!black] (0.15,0.15) -- (0.15,1.65)
      node[above] {threshold branch point};
    \node[align=left] at (3.2,1.3)
      {physical sheet above the cut\\second sheet below the cut};
  \end{tikzpicture}
  \caption{Schematic energy-plane geometry for a single threshold at $E=0$.
  Complex scaling by an angle $\theta>0$ rotates the continuous spectrum by
  $-2\theta$ and exposes a resonance pole.  The drawing is topological: the
  precise sheet structure and the number of thresholds depend on the problem.}
  \label{fig:unified-plane}
\end{figure}

\section{Scope and terminology}
\label{subsec:scope}
% BEGIN CURATED SUBJECT INDEX
\index{operator!adjoint and self-adjoint}
\index{unbounded operator}
\index{Hilbert space}
\index{spectrum!point}
\index{resolvent}
\index{Riemann sheet}
\index{resonance!residue}
\index{antiresonance}
\index{Breit--Wigner approximation}
% END CURATED SUBJECT INDEX

For a bounded finite-dimensional matrix, ``Hermitian'' and ``self-adjoint''
are interchangeable.  For a differential operator they are not.  Symmetry of
the differential expression does not fix the operator domain, and boundary
conditions can change the adjoint.  We therefore use \textit{self-adjoint} for
operators on a Hilbert space and reserve \textit{Hermitian} for matrices or
informal descriptions.  The complex-scaled Hamiltonian is generally
non-self-adjoint even when the original Hamiltonian is self-adjoint.

We use \textit{resonance} primarily for a pole of the continued resolvent or
$S$ matrix.  The word \textit{resonant state} denotes the residue vector or
functional associated with that pole.  A Breit--Wigner peak is evidence for a
nearby simple pole, but background interference, thresholds, and overlapping
poles can obscure or even eliminate a visible peak.  Conversely, a complex
eigenvalue of an arbitrary non-Hermitian truncation need not be a resonance;
stability under changes of deformation angle, basis, and truncation is
essential.

We distinguish resonance poles from three neighboring notions.  A bound-state
pole lies on the physical sheet below threshold and has a square-integrable
eigenfunction.  A virtual or antibound state lies on a nonphysical sheet on an
imaginary momentum axis and need not produce a narrow peak.  An antiresonance
is the time-reversed partner of a resonance, related by the appropriate
reality symmetry of the potential.  These classifications are simplest in the
complex momentum plane and become more intricate for multiple channels.

The main construction concerns stationary spectral theory.  A resonance pole
produces exponential decay only over an intermediate time interval.  The
short-time survival probability is quadratic, while thresholds generate
long-time power-law tails
\cite{GarciaCalderon2017NonExp,Giacosa2012NonExpQFTQM,
Garmon2017NonExpThreshold}.  These deviations are not failures of the pole
description; they are contributions from the analytic background and branch
cuts that accompany the pole.

\section{Organization and the central line of argument}
\label{subsec:organization}
% BEGIN CURATED SUBJECT INDEX
\index{resolvent}
\index{complex scaling}
\index{CDCC}
\index{lithium nuclei}
\index{beryllium nuclei}
\index{exceptional point}
\index{horizon boundary condition}
% END CURATED SUBJECT INDEX

Section~\ref{sec:computational-spine} states the algorithm used throughout.
Eight Stages then add one obstruction at a time.  Stage I constructs the
spectral and operator-algebraic language; Stage II derives scattering data
and tests its short-range hypotheses.  Stage III builds the exact-WKB global
connection datum, while Stage IV proves how the same datum is represented by
resolvents, regularized states, complex scaling, rigged spaces, and continuum
response.

Stage V adds a radial singular endpoint.  Stage VI enlarges the projectile to
a weakly bound few-body system and derives CDCC, including lithium and
beryllium-core reaction examples.  Stage VII treats an exceptional point as a
multiple zero of the analytic pole condition.  Stage VIII derives black-hole
master equations and turns horizon boundary conditions into the same
connection problem.  The final Part separates invariant data from contour,
basis, regulator, and representation choices.

Detailed calculations developed in Refs.~\cite{Morikawa:2025grx,
Morikawa:2025xjq,Morikawa:2025vvs,Morikawa:2025ezl,Morikawa:2025inx,
Ogawa:2026veu,Ogawa:2026dSCSM} supply seven of the Laboratories.  They have
been reorganized rather than reproduced as papers: common scattering,
exact-WKB, complex-scaling, and black-hole preliminaries are established once
in canonical chapters, and each Laboratory begins at the first calculation
specific to its model.  The remaining two Laboratories develop complementary
few-body nuclear branches.  Every one is meant to be read with pencil, paper,
and a numerical notebook; its closing checkpoint states what must be reproducible before the
reader proceeds.
